\clearpage
\begin{examplebox}
\textbf{\textcolor{CExample}{例 1｜从一张数字表看懂“近小远大”（最基础）}}
\begin{center}\begin{tikzpicture}[font=\small,>=stealth,line join=round]
\draw[gray!40,rounded corners=3pt] (0,0) rectangle (14.2,8.3);\node at (7.1,7.700000000000001) {只看允许段：最近处最低，最远处最高};
\begin{scope}[shift={(1.8777777777777778,2.2249999999999996)},x=1.111111111111111cm,y=0.575cm]\draw[->,gray] (-0.7,0)--(4.7,0) node[right]{$x$};\draw[->,gray] (0,-1)--(0,7) node[above]{$y$};\draw[gray,thick,solid,domain=0:4,samples=180] plot(\x,{(\x-2)^2+1});\draw[blue,thick,very thick,domain=0:4,samples=180] plot(\x,{(\x-2)^2+1});\fill[orange!85!black] (0,5) circle(2.2pt);\fill[orange!85!black] (4,5) circle(2.2pt);\fill[orange!85!black] (2,1) circle(2.2pt);\draw[violet,dashed] (2,-.2)--(2,6.5);\draw[orange!85!black,very thick] (0,-.5)--(4,-.5);\node[above left] at (0,5){$(0,5)$};\node[above right] at (4,5){$(4,5)$};\node[below right] at (2,1){$(2,1)$};\end{scope}\node[anchor=west,align=left,text width=5.0cm] at (8.4,4.0) {$f=(x-2)^2+1$\\[9pt]$x\in[0,4]$\\[9pt]最近：$x=2$，最小值1\\[9pt]最远：$x=0,4$\\[9pt]最大值5，两端都取到};
\node[text width=12.6cm,align=center] at (7.1,.65) {图取a为正；开口向下时，最近与最远对应的最大最小关系反过来。};
\end{tikzpicture}\end{center}



\textbf{题目：} $f(x)=(x-2)^2+1$，只让 $x\in[0,4]$，求最大值、最小值及取到的位置。

\textbf{第一步：} 对称轴 $x=2$，位于 $[0,4]$ 内；开口向上。

\textbf{第二步：} 离 $2$ 最近的地方是 $x=2$，故最小值 $f(2)=1$。离 $2$ 最远的是 $x=0$ 和 $x=4$（都是相距 $2$），故最大值 $f(0)=f(4)=5$。

\textbf{答案：} $\boxed{f_{\min}=1\ (x=2),\quad f_{\max}=5\ (x=0\text{ 或 }4)}$。

\textbf{看懂了什么？} 两个端点可以\textbf{同时取到同一个最值}，不是只能选一个。

\medskip

\end{examplebox}
\clearpage
\begin{examplebox}
\textbf{\textcolor{CExample}{例 2｜原题基础：顶点在内部，开口向上}}
\begin{center}\begin{tikzpicture}[font=\small,>=stealth,line join=round]
\draw[gray!40,rounded corners=3pt] (0,0) rectangle (14.2,8.3);\node at (7.1,7.700000000000001) {例2：顶点在中点，两端同时达到最大值};
\begin{scope}[shift={(2.988888888888889,2.7115384615384612)},x=1.111111111111111cm,y=0.3538461538461538cm]\draw[->,gray] (-1.7,0)--(3.7,0) node[right]{$x$};\draw[->,gray] (0,-3)--(0,10) node[above]{$y$};\draw[blue,thick,very thick,domain=-1:3,samples=180] plot(\x,{2*(\x-1)^2-1});\fill[orange!85!black] (-1,7) circle(2.2pt);\fill[orange!85!black] (1,-1) circle(2.2pt);\fill[orange!85!black] (3,7) circle(2.2pt);\draw[violet,dashed] (1,-.8)--(1,9);\node[above] at (-1,7){$(-1,7)$};\node[above] at (3,7){$(3,7)$};\node[below] at (1,-1){$(1,-1)$};\end{scope}\node[anchor=west,align=left,text width=5.0cm] at (8.4,4.0) {$f=2(x-1)^2-1$\\[9pt]$x\in[-1,3]$\\[9pt]近处：$x=1$，最小$-1$\\[9pt]两端到轴距离都为2\\[9pt]最大值 $f(-1)=f(3)=7$};
\node[text width=12.6cm,align=center] at (7.1,.65) {同一最大值可以在两个不同位置取得；这里两个端点关于轴对称。};
\end{tikzpicture}\end{center}



\textbf{题目：} $f(x)=2x^2-4x+1$，$x\in[-1,3]$，求最大值与最小值。

\textbf{第一步（配方）：}
\[
 f(x)=2(x-1)^2-1.
\]
所以对称轴 $x=1$，顶点的值为 $-1$，开口向上。

\textbf{第二步（比较远近）：} $1$ 在区间内，最近处就是 $x=1$；左端 $-1$ 与 $1$ 相差 $2$，右端 $3$ 与 $1$ 也相差 $2$，所以两端同高。

\textbf{第三步（计算）：}
\[
 f(1)=-1,\qquad f(-1)=f(3)=7.
\]
\textbf{答案：} $\boxed{f_{\min}=-1,\qquad f_{\max}=7}$；最大值在两个端点都能取到。

\medskip

\end{examplebox}
\clearpage
\begin{examplebox}
\textbf{\textcolor{CExample}{例 3｜原题稍变形：开口向下，最远反而最小}}
\begin{center}\begin{tikzpicture}[font=\small,>=stealth,line join=round]
\draw[gray!40,rounded corners=3pt] (0,0) rectangle (14.2,8.3);\node at (7.1,7.700000000000001) {例3：开口向下，最远的右端取得最小值};
\begin{scope}[shift={(2.0545454545454547,2.8)},x=1.3636363636363635cm,y=0.575cm]\draw[->,gray] (-0.7,0)--(3.7,0) node[right]{$x$};\draw[->,gray] (0,-2)--(0,6) node[above]{$y$};\draw[blue,thick,very thick,domain=0:3,samples=180] plot(\x,{-(\x-1)^2+4});\fill[orange!85!black] (0,3) circle(2.2pt);\fill[orange!85!black] (1,4) circle(2.2pt);\fill[orange!85!black] (3,0) circle(2.2pt);\node[left] at (0,3){$(0,3)$};\node[above] at (1,4){$(1,4)$};\node[below] at (3,0){$(3,0)$};\end{scope}\node[anchor=west,align=left,text width=5.0cm] at (8.4,4.0) {$f=-(x-1)^2+4$\\[9pt]$x\in[0,3]$\\[9pt]轴 $x=1$\\[9pt]最近取最大值4\\[9pt]最远$x=3$，最小值0};
\node[text width=12.6cm,align=center] at (7.1,.65) {左端距离1、右端距离2；更远的端点反而更低。};
\end{tikzpicture}\end{center}



\textbf{题目：} $f(x)=-x^2+2x+3$，$x\in[0,3]$，求最大值、最小值。

\textbf{第一步：} $f(x)=-(x-1)^2+4$，对称轴 $x=1$ 在区间内，开口向下。

\textbf{第二步：} 离 $1$ 最近是 $x=1$，故最大值为 $f(1)=4$。两个端点离对称轴的距离分别为 $1,2$，较远的是 $x=3$，故最小值为 $f(3)=0$。

\textbf{答案：} $\boxed{f_{\max}=4,\quad f_{\min}=0}$。

\textbf{容易错在哪？} 开口向下时，不是“越远越大”，而是\textbf{越远越小}。

\medskip

\end{examplebox}
\clearpage
\begin{examplebox}
\textbf{\textcolor{CExample}{例 4｜对称轴不在区间内：为什么不能用顶点值？}}
\begin{center}\begin{tikzpicture}[font=\small,>=stealth,line join=round]
\draw[gray!40,rounded corners=3pt] (0,0) rectangle (14.2,8.3);\node at (7.1,7.700000000000001) {顶点虽低，却不在允许的区间中};
\begin{scope}[shift={(1.8777777777777778,2.2249999999999996)},x=1.111111111111111cm,y=0.575cm]\draw[->,gray] (-0.7,0)--(4.7,0) node[right]{$x$};\draw[->,gray] (0,-1)--(0,7) node[above]{$y$};\draw[gray,thick,solid,domain=-0.2:4.2,samples=180] plot(\x,{(\x-2)^2+1});\draw[blue,thick,very thick,domain=3:4,samples=180] plot(\x,{(\x-2)^2+1});\fill[orange!85!black] (3,2) circle(2.2pt);\fill[orange!85!black] (4,5) circle(2.2pt);\draw[violet,dashed] (2,1.05)--(2,6.5);\draw[orange!85!black,very thick] (3,-.5)--(4,-.5);\node[below] at (2,1){$(2,1)$};\node[right] at (3,2){$(3,2)$};\node[above left=3pt] at (4,5){$(4,5)$};\foreach \v in {2,3,4}{\node[below=5pt] at (\v,-.5){$\v$};}\end{scope}\node[anchor=west,align=left,text width=5.0cm] at (8.4,4.0) {$f=(x-2)^2+1$\\[9pt]允许 $x\in[3,4]$\\[9pt]$2\notin[3,4]$\\[9pt]最近 $x=3$：最小值2\\[9pt]最远 $x=4$：最大值5};
\node[text width=12.6cm,align=center] at (7.1,.65) {灰线仅供定位；蓝色粗线是允许段，不能到灰色顶点处取值。};
\end{tikzpicture}\end{center}



\textbf{题目：} $f(x)=(x-2)^2+1$，$x\in[3,4]$，求最值。

\textbf{第一步：} 顶点在 $x=2$，但题目只允许 $3\le x\le4$。

\textbf{第二步：} 区间全在对称轴右侧，越往右距离越大；开口向上，所以越往右函数值越大。

\textbf{第三步：}
\[
 f(3)=(3-2)^2+1=2,\qquad f(4)=5.
\]
\textbf{答案：} $\boxed{f_{\min}=2,\quad f_{\max}=5}$。

\textbf{警醒：} 虽然整条抛物线最低可以到 $1$，但在指定区间内\textbf{最低只能到 $2$}。

\medskip

\end{examplebox}
\clearpage
\begin{examplebox}
\textbf{\textcolor{CExample}{例 5｜计算加速：区间中点一眼判断端点}}
\begin{center}\begin{tikzpicture}[font=\small,>=stealth,line join=round]
\draw[gray!40,rounded corners=3pt] (0,0) rectangle (14.2,8.3);\node at (7.1,7.700000000000001) {例5：轴偏左，右端远，直接选出最大值位置};
\begin{scope}[shift={(1.6675675675675676,2.3576923076923073)},x=0.8108108108108107cm,y=0.044230769230769226cm]\draw[->,gray] (-0.7,0)--(6.7,0) node[right]{$x$};\draw[->,gray] (0,-16)--(0,88) node[above]{$y$};\draw[blue,thick,very thick,domain=0:6,samples=180] plot(\x,{3*(\x-1)^2-2});\fill[orange!85!black] (0,1) circle(2.2pt);\fill[orange!85!black] (1,-2) circle(2.2pt);\fill[orange!85!black] (6,73) circle(2.2pt);\draw[violet,dashed] (1,-1)--(1,82);\draw[gray,dashed] (3,0)--(3,82);\node[above] at (1,82){$h=1$};\node[above] at (3,82){$s=3$};\node[above left] at (6,73){$(6,73)$};\node[below right] at (1,-2){$(1,-2)$};\foreach \v in {0,1,3,6}{\node[below=3pt] at (\v,-12){$\v$};}\end{scope}\node[anchor=west,align=left,text width=5.0cm] at (8.4,4.0) {$f=3(x-1)^2-2$\\[9pt]区间 $[0,6]$，中点3\\[9pt]$h=1<3$，右端更远\\[9pt]$f(6)=73$ 为最大值\\[9pt]$f(1)=-2$ 为最小值};
\node[text width=12.6cm,align=center] at (7.1,.65) {纵轴使用原函数值；近处曲线显平是数值跨度大，坐标标注给出准确最值。};
\end{tikzpicture}\end{center}



\textbf{题目：} $f(x)=3(x-1)^2-2$，$x\in[0,6]$，不分别计算两个端点，直接指出最大值在哪个端点，并求最大最小值。

\textbf{第一步：} 对称轴 $h=1$，区间中点 $s=3$，所以对称轴\textbf{在中点左边}。右端点 $x=6$ 离对称轴更远。

\textbf{第二步：} 开口向上，故最大值在 $x=6$；$h=1$ 又在区间内，最小值在 $x=1$。
\[
 f_{\max}=f(6)=3\times25-2=73,\qquad f_{\min}=f(1)=-2.
\]
\textbf{答案：} $\boxed{f_{\max}=73,\quad f_{\min}=-2}$。

\textbf{速度提升：} 整道题只需要算两次值：$f(6)$ 和 $f(1)$，不用先算 $f(0)$ 再比较。

\medskip

\end{examplebox}
\clearpage
\begin{examplebox}
\textbf{\textcolor{CExample}{例 6｜反过来求参数：最小值必须真的取到}}
\begin{center}\begin{tikzpicture}[font=\small,>=stealth,line join=round]
\draw[gray!40,rounded corners=3pt] (0,0) rectangle (14.2,8.5);\node at (7.1,7.9) {例6：要取到最小值2，区间必须包含1};
\node at (2.4,6.9) {$m=\frac12$};\begin{scope}[shift={(1.2666666666666666,3.0)},x=1.1333333333333333cm,y=0.8cm]\draw[->,gray] (-0.5,0)--(2.5,0) node[right]{$x$};\draw[->,gray] (0,-0.5)--(0,4) node[above]{$y$};\draw[gray,thick,solid,domain=-0.1:2.1,samples=160] plot(\x,{(\x-1)^2+2});\draw[blue,thick,very thick,domain=0:0.5,samples=160] plot(\x,{(\x-1)^2+2});\fill[orange!85!black] (0.5,2.25) circle(2pt);\fill[orange!85!black] (0,3) circle(2pt);\end{scope}\node[align=center,text width=4cm] at (2.4,1.7) {取不到2};\node at (7.050000000000001,6.9) {$m=1$};\begin{scope}[shift={(5.916666666666667,3.0)},x=1.1333333333333333cm,y=0.8cm]\draw[->,gray] (-0.5,0)--(2.5,0) node[right]{$x$};\draw[->,gray] (0,-0.5)--(0,4) node[above]{$y$};\draw[gray,thick,solid,domain=-0.1:2.1,samples=160] plot(\x,{(\x-1)^2+2});\draw[blue,thick,very thick,domain=0:1,samples=160] plot(\x,{(\x-1)^2+2});\fill[orange!85!black] (1,2) circle(2pt);\fill[orange!85!black] (0,3) circle(2pt);\fill[orange!85!black] (1,2) circle(2pt);\end{scope}\node[align=center,text width=4cm] at (7.050000000000001,1.7) {可在$x=1$取到2};\node at (11.7,6.9) {$m=2$};\begin{scope}[shift={(10.566666666666666,3.0)},x=1.1333333333333333cm,y=0.8cm]\draw[->,gray] (-0.5,0)--(2.5,0) node[right]{$x$};\draw[->,gray] (0,-0.5)--(0,4) node[above]{$y$};\draw[gray,thick,solid,domain=-0.1:2.1,samples=160] plot(\x,{(\x-1)^2+2});\draw[blue,thick,very thick,domain=0:2,samples=160] plot(\x,{(\x-1)^2+2});\fill[orange!85!black] (2,3) circle(2pt);\fill[orange!85!black] (0,3) circle(2pt);\fill[orange!85!black] (1,2) circle(2pt);\end{scope}\node[align=center,text width=4cm] at (11.7,1.7) {可在$x=1$取到2};
\node[text width=12.6cm,align=center] at (7.1,.65) {同尺度，$f=(x-1)^2+2$；蓝色为$[0,m]$，恰当条件是$m\ge1$。};
\end{tikzpicture}\end{center}



\textbf{题目（保留原题并讲清条件）：} 已知 $f(x)=x^2-2x+3$ 在闭区间 $[0,m]$ 上的最小值是 $2$。约定 $m>0$，求 $m$ 的范围。

\textbf{第一步：} 配方得
\[
 f(x)=(x-1)^2+2.
\]
因为平方永远不小于 $0$，所以不管选什么 $x$，$f(x)\ge2$。\textbf{要真正等于 $2$，平方必须等于 $0$，也就是 $x=1$。}

\textbf{第二步：} 题目只允许 $x\in[0,m]$，所以必须让 $1$ 在区间里：
\[
 0\le1\le m \quad\Longrightarrow\quad m\ge1.
\]
\textbf{答案：} $\boxed{m\in[1,+\infty)}$。

\textbf{为什么比直接分类更快？} 不用先把 $m<1$、$m\ge1$ 两种情况全部算完，只需想清楚\textbf{最小值 $2$ 只能在 $x=1$ 取得}。\textbf{提醒：} 本题的区间右端点在变，是对固定区间结论的进一步应用。

\medskip

\end{examplebox}
\clearpage
\begin{examplebox}
\textbf{\textcolor{CExample}{例 7｜带参数的对称轴：分界点为什么是中点？（拓展）}}
\begin{center}\begin{tikzpicture}[font=\small,>=stealth,line join=round]
\draw[gray!40,rounded corners=3pt] (0,0) rectangle (14.2,8.5);\node at (7.1,7.9) {例7：最近点看0和4，最远端点看中点2};
\node at (2.4,6.9) {$t=-1$};\begin{scope}[shift={(1.04,2.9272727272727272)},x=0.6799999999999999cm,y=0.1090909090909091cm]\draw[->,gray] (-0.5,0)--(4.5,0) node[right]{$x$};\draw[->,gray] (0,-3)--(0,30) node[above]{$y$};\draw[blue,thick,very thick,domain=0:4,samples=160] plot(\x,{(\x-(-1))^2+1});\fill[orange!85!black] (0,2) circle(2pt);\fill[orange!85!black] (4,26) circle(2pt);\end{scope}\node[align=center,text width=4cm] at (2.4,1.7) {最近0，最小2\\最远4，最大26};\node at (7.050000000000001,6.9) {$t=2$};\begin{scope}[shift={(5.69,2.9272727272727272)},x=0.6799999999999999cm,y=0.1090909090909091cm]\draw[->,gray] (-0.5,0)--(4.5,0) node[right]{$x$};\draw[->,gray] (0,-3)--(0,30) node[above]{$y$};\draw[blue,thick,very thick,domain=0:4,samples=160] plot(\x,{(\x-(2))^2+1});\fill[orange!85!black] (0,5) circle(2pt);\fill[orange!85!black] (4,5) circle(2pt);\fill[orange!85!black] (2,1) circle(2pt);\end{scope}\node[align=center,text width=4cm] at (7.050000000000001,1.7) {最近2，最小1\\两端同远，最大5};\node at (11.7,6.9) {$t=5$};\begin{scope}[shift={(10.34,2.9272727272727272)},x=0.6799999999999999cm,y=0.1090909090909091cm]\draw[->,gray] (-0.5,0)--(4.5,0) node[right]{$x$};\draw[->,gray] (0,-3)--(0,30) node[above]{$y$};\draw[blue,thick,very thick,domain=0:4,samples=160] plot(\x,{(\x-(5))^2+1});\fill[orange!85!black] (0,26) circle(2pt);\fill[orange!85!black] (4,2) circle(2pt);\end{scope}\node[align=center,text width=4cm] at (11.7,1.7) {最近4，最小2\\最远0，最大26};
\node[text width=12.6cm,align=center] at (7.1,.65) {三图同尺度；区间固定为$[0,4]$，对称轴$x=t$改变。};
\end{tikzpicture}\end{center}



\textbf{题目：} $g(x)=(x-t)^2+1$，$x\in[0,4]$，其中 $t$ 是实数。求最小值与最大值。

\textbf{最小值：} 开口向上，找离 $x=t$ 最近的位置：
\[
 g_{\min}=\begin{cases}
 t^2+1,&t<0\quad(x=0),\\
 1,&0\le t\le4\quad(x=t),\\
 (4-t)^2+1,&t>4\quad(x=4).
 \end{cases}
\]
\textbf{最大值：} 开口向上，找最远端点。区间中点为 $2$，所以分界点\textbf{不是 $0$ 或 $4$，而是 $2$}：
\[
 g_{\max}=\begin{cases}
 (4-t)^2+1,&t<2\quad(x=4),\\
 5,&t=2\quad(x=0\text{ 或 }4),\\
 t^2+1,&t>2\quad(x=0).
 \end{cases}
\]
\textbf{看出什么？} 讨论最小值，关键是对称轴是否进入区间（$0,4$）；讨论最大值，关键是对称轴跨过\textbf{区间中点 $2$}。这是连接“轴动区间定”题型的一把钥匙。
\end{examplebox}
